\begin{tikzpicture}[font=\small]
  % ---------- DT rectangular pulse ----------
  \begin{scope}
    \draw[ax] (-3.6,0) -- (3.8,0) node[below left]{$n$};
    \draw[ax] (0,-0.2) -- (0,1.9);
    \node[note, anchor=south east] at (-0.1,1.5){$x[n]$};
    \foreach \n in {-2,-1,0,1,2} \stemat{ssblue}{\n*0.62}{1.3}
    \foreach \n in {-5,-4,-3,3,4,5} \filldraw[sslight] (\n*0.62,0) circle (1.7pt);
    \node[note, anchor=north] at (-1.24,-0.12){$-N_1$};
    \node[note, anchor=north] at (1.24,-0.12){$N_1$};
    \node[note, anchor=north, align=center] at (0,-0.75){矩形序列（$|n|\le N_1$ 取 1）};
  \end{scope}

  % ---------- its DTFT: periodic Dirichlet ----------
  \begin{scope}[xshift=11.6cm]
    \draw[ax] (-6.6,0) -- (6.9,0) node[below left]{$\Omega$};
    \draw[ax] (0,-0.9) -- (0,2.3);
    \node[note, anchor=south east] at (-0.15,1.35){$X(e^{j\Omega})$};

    % Dirichlet kernel, period 2*pi mapped to 4.0 units
    \foreach \s in {-1,0,1}{
      \draw[curve, domain=-1.98:1.98, samples=300, smooth, shift={(\s*4.0,0)}]
        plot (\x, {2.0*sin(deg(5*(\x*3.14159/4.0)+0.00001))/(5*sin(deg((\x*3.14159/4.0)+0.00001))+0.00001)*0.5});
    }
    \draw[helper] (2.0,-0.9) -- (2.0,2.1);
    \draw[helper] (-2.0,-0.9) -- (-2.0,2.1);
    \draw[helper] (6.0,-0.9) -- (6.0,2.1);
    \draw[helper] (-6.0,-0.9) -- (-6.0,2.1);
    \node[note, anchor=north] at (2.0,-0.95){$\pi$};
    \node[note, anchor=north] at (-2.0,-0.95){$-\pi$};
    \node[note, anchor=north] at (4.0,-0.95){$2\pi$};
    \node[note, anchor=north] at (-4.0,-0.95){$-2\pi$};

    \draw[{Stealth[length=4pt]}-{Stealth[length=4pt]}, ssred] (-2.0,2.25) -- (2.0,2.25);
    \node[ssred, font=\footnotesize, anchor=south] at (0,2.28){一个周期 $2\pi$};

    \node[note, anchor=north, align=center] at (0,-1.55)
      {$X(e^{j\Omega})=\dfrac{\sin\!\big(\Omega(N_1+\tfrac12)\big)}{\sin(\Omega/2)}$};
  \end{scope}

  \node[note, anchor=north, align=center] at (5.6,-3.4)
    {DTFT 与 CTFT 最大的差别：它对 $\Omega$ 恒为 $2\pi$ 周期。\;
     因为 $e^{j(\Omega+2\pi)n}=e^{j\Omega n}$ —— 离散时间里高于 $\pi$ 的频率根本不存在，\\
     $\Omega=\pi$ 就是「最高频率」（相邻样点正负交替），$\Omega=0$ 与 $2\pi$ 都是直流};
\end{tikzpicture}
